\chapter{Reciprocal Compactification and the Prime-Tail Certificate}
\label{ch:prime-tail-certificate}

The Hermite ladder of \Cref{ch:phasenav-weil-hermite-ladder} replaces one
localized witness by a concrete dense family and finite principal sections.
Its arithmetic matrices are nevertheless evaluated with a finite prime-power
cutoff. This chapter supplies an explicit error certificate for removing that
cutoff in every fixed finite Hermite section.

The reciprocal map used here is not a map of zeta zeros. It acts only on the
logarithmic coordinate of the omitted prime tail:
\[
 u=\log x,
 \qquad
 z_{\rm t}=\frac1u.
\]
Its purpose is analytic compactification: the half-line
$u\in[\log Q,\infty)$ becomes the finite interval
$z_{\rm t}\in(0,1/\log Q]$.

\section{The omitted prime-power matrix}

Fix the Gaussian width $w>0$ and the translated Hermite channels from
\Cref{ch:phasenav-weil-hermite-ladder}. Under the Fourier convention
\[
 \widehat f(\xi)
 =
 \int_{\R} f(r)e^{-2\pi i\xi r}\,dr,
\]
the product kernel has the form
\[
 \widehat H_{mn}(\xi)
 =
 A_{mn}\,
 e^{-2\pi i\gamma_0\xi}
 e^{-\kappa^2/4}
 \sum_{d\in D_{mn}} c_{mn,d}(-i\kappa)^d,
 \qquad
 \kappa=\frac{2\pi\xi}{w},
\]
where
\[
 A_{mn}
 =
 \frac{\sqrt{\pi}}{w}
 \left(
 \frac{w}{\sqrt{\pi}\,2^m m!}
 \right)^{1/2}
 \left(
 \frac{w}{\sqrt{\pi}\,2^n n!}
 \right)^{1/2},
\]
and the degrees and coefficients are obtained from the exact Hermite
linearization
\[
 H_m(y)H_n(y)
 =
 \sum_{k=0}^{\min(m,n)}
 2^k k!\binom{m}{k}\binom{n}{k}
 H_{m+n-2k}(y).
\]

For a cutoff $Q\ge3$, define the omitted prime-power contribution
\[
 E_{mn}(Q)
 =
 -\frac1{2\pi}
 \sum_{\substack{q>Q\\q=p^\ell}}
 \frac{\Lambda(q)}{\sqrt q}
 \left[
 \widehat H_{mn}\!\left(\frac{\log q}{2\pi}\right)
 +
 \widehat H_{mn}\!\left(-\frac{\log q}{2\pi}\right)
 \right].
\]
No zero list occurs in this definition.

\section{Monotonicity beyond an explicit threshold}

For a polynomial degree $d\ge0$, introduce
\[
 g_d(x)
 =
 x^{-1/2}(\log x)^{d+1}
 \exp\!\left[-\frac{(\log x)^2}{4w^2}\right],
 \qquad x>1.
\]

\begin{lemma}[Explicit monotonicity threshold]
\label{lem:tail-monotonicity}
The function $g_d$ is decreasing whenever
\[
 \log x\ge
 \tau_d(w)
 :=
 \frac{\sqrt{w^4+8w^2(d+1)}-w^2}{2}.
\]
\end{lemma}

\begin{proof}
Put $u=\log x$. The logarithmic derivative with respect to $u$ is
\[
 \frac{d}{du}\log g_d(e^u)
 =
 -\frac12+\frac{d+1}{u}-\frac{u}{2w^2}.
\]
It is non-positive precisely when
\[
 u^2+w^2u-2w^2(d+1)\ge0.
\]
The positive root of the quadratic is $\tau_d(w)$.
\end{proof}

Consequently, when $\log Q\ge\tau_d(w)$, the integral test gives
\[
 \sum_{n>Q}g_d(n)
 \le
 \int_Q^\infty g_d(x)\,dx.
\]
This step uses only positivity and monotonicity. It does not use the prime
number theorem or any hypothesis about zeta zeros.

\section{Closed incomplete-gamma form}

Define
\[
 I_d(Q,w)
 =
 \frac1{w^d}
 \int_Q^\infty
 (\log x)^{d+1}x^{-1/2}
 \exp\!\left[-\frac{(\log x)^2}{4w^2}\right]dx.
\]

\begin{theorem}[Prime-tail integral in closed form]
\label{thm:tail-gamma}
Let
\[
 a_Q=\frac{\log Q-w^2}{2w}.
\]
For $\log Q>w^2$,
\[
 I_d(Q,w)
 =
 \frac{e^{w^2/4}}{w^d}
 \sum_{j=0}^{d+1}
 \binom{d+1}{j}
 w^{2(d+1-j)}
 2^j w^{j+1}
 \Gamma\!\left(\frac{j+1}{2},a_Q^2\right),
\]
where $\Gamma(s,a)$ is the upper incomplete gamma function.
\end{theorem}

\begin{proof}
With $u=\log x$,
\[
 I_d(Q,w)
 =
 \frac1{w^d}
 \int_{\log Q}^{\infty}
 u^{d+1}
 \exp\!\left(-\frac{u^2}{4w^2}+\frac u2\right)du.
\]
Complete the square:
\[
 -\frac{u^2}{4w^2}+\frac u2
 =
 \frac{w^2}{4}
 -
 \frac{(u-w^2)^2}{4w^2}.
\]
Set $a=(u-w^2)/(2w)$, expand
$(2wa+w^2)^{d+1}$, and use
\[
 \int_{a_Q}^{\infty}a^j e^{-a^2}\,da
 =
 \frac12\Gamma\!\left(\frac{j+1}{2},a_Q^2\right).
\]
Collecting powers gives the stated expression.
\end{proof}

\section{Reciprocal compactification}

The same integral admits a finite-interval representation. Set
\[
 z_{\rm t}=\frac1u=\frac1{\log x}.
\]
Then
\[
 I_d(Q,w)
 =
 \frac1{w^d}
 \int_0^{1/\log Q}
 z_{\rm t}^{-(d+3)}
 \exp\!\left[
 -\frac1{4w^2z_{\rm t}^2}
 +\frac1{2z_{\rm t}}
 \right]dz_{\rm t}.
\]

\begin{theorem}[Flat reciprocal endpoint]
\label{thm:tail-reciprocal-flat}
Define the integrand at $z_{\rm t}=0$ to be zero. The reciprocal tail
integrand extends smoothly to the closed interval
$[0,1/\log Q]$, and every derivative at the new endpoint vanishes.
\end{theorem}

\begin{proof}
For sufficiently small positive $z_{\rm t}$, the exponent is dominated by
$-1/(4w^2z_{\rm t}^2)$. Multiplication by any fixed negative power of
$z_{\rm t}$, including those produced by repeated differentiation, cannot
overcome this Gaussian decay. Hence the function and all of its derivatives
tend to zero at the endpoint.
\end{proof}

The map $z_{\rm t}=1/\log x$ therefore converts the infinite tail into a
compact quadrature domain without introducing a singular boundary
contribution. This is the precise role of the reciprocal map.

\section{Entrywise prime-tail certificate}

The elementary von Mangoldt bound
\[
 0\le\Lambda(n)\le\log n
\]
and the triangle inequality imply the following unconditional estimate.

\begin{theorem}[Fixed-entry cutoff certificate]
\label{thm:entry-tail-certificate}
Assume
\[
 \log Q\ge
 \max_{d\in D_{mn}}\tau_d(w).
\]
Then
\[
 |E_{mn}(Q)|
 \le
 B_{mn}(Q)
 :=
 \frac{A_{mn}}{\pi}
 \sum_{d\in D_{mn}}
 c_{mn,d}\,I_d(Q,w).
\]
In particular, $B_{mn}(Q)\to0$ as $Q\to\infty$ for each fixed pair
$(m,n)$.
\end{theorem}

\begin{proof}
At the prime-power point $\xi=\log q/(2\pi)$ one has
$\kappa=\log q/w$. The unit-modulus translation phase disappears after
taking absolute values, while the two Fourier terms contribute a factor at
most two. Thus
\[
 |E_{mn}(Q)|
 \le
 \frac{A_{mn}}{\pi}
 \sum_{d\in D_{mn}}c_{mn,d}
 \sum_{\substack{q>Q\\q=p^\ell}}
 \frac{\Lambda(q)}{\sqrt q}
 \left(\frac{\log q}{w}\right)^d
 e^{-(\log q)^2/(4w^2)}.
\]
Replace the prime-power support by all integers, use
$\Lambda(q)\le\log q$, and apply \Cref{lem:tail-monotonicity} and the
integral test. The result is exactly the stated sum of $I_d(Q,w)$.
\end{proof}

\section{Finite-section norm and eigenvalue enclosure}

Let $E_N(Q)$ be the $N\times N$ Hermitian matrix with entries $E_{mn}(Q)$,
and let $B_N(Q)$ be the non-negative symmetric matrix with entries
$B_{mn}(Q)$ for $0\le m,n<N$.

\begin{corollary}[Operator-norm certificate]
\label{cor:operator-tail-certificate}
For every fixed finite section,
\[
 \|E_N(Q)\|_2
 \le
 \|B_N(Q)\|_\infty
 =
 \max_{0\le m<N}
 \sum_{n=0}^{N-1}B_{mn}(Q).
\]
If $W_N(Q)$ is the cutoff matrix and $W_N(\infty)$ its convergent
prime-side limit, then each ordered eigenvalue satisfies
\[
 \left|
 \lambda_j\!\left(W_N(\infty)\right)
 -
 \lambda_j\!\left(W_N(Q)\right)
 \right|
 \le
 \|B_N(Q)\|_\infty.
\]
\end{corollary}

\begin{proof}
Entrywise domination gives
$\|E_N(Q)\|_\infty\le\|B_N(Q)\|_\infty$. Since $E_N(Q)$ is Hermitian,
its spectral norm is bounded by the maximum absolute row sum. The
eigenvalue statement is Weyl's perturbation inequality.
\end{proof}

This replaces an empirical difference between two cutoffs by a one-sided
certificate for the entire omitted tail.

\section{Declared numerical receipt}

The native programme
\[
 \texttt{construction/phasenav/secret\_of\_half\_weil\_prime\_tail\_certificate.pnv}
\]
declares
\[
 w=0.8,
 \qquad
 Q=100000,
 \qquad
 1\le N\le6.
\]
The executor evaluates the gamma expression, direct logarithmic integral, and
reciprocal compact-interval integral independently. Their largest relative
disagreement is below the declared high-precision tolerance.

For the largest declared section,
\[
 \boxed{
 \|E_6(100000)\|_2
 \le
 7.717202888999335\times10^{-13}
 }.
\]
Finite shells between $Q$ and $2Q$ are also checked directly against the full
tail majorant. These shell checks are regression tests, not ingredients of
the analytic proof.

\section{Claim boundary}

The results of this chapter establish:

\begin{description}
\item[SOH-L018, exact.]
The reciprocal logarithmic coordinate compactifies the prime tail to a finite
interval with a smooth flat endpoint.

\item[SOH-L019, exact.]
Every fixed Hermite matrix entry has an explicit incomplete-gamma majorant
under the stated monotonicity condition.

\item[SOH-L020, exact finite-section theorem.]
The entrywise certificates produce an operator-norm error bound and an
eigenvalue enclosure for every fixed finite principal section.

\item[SOH-N005, numerical certificate.]
For the declared $w$, $Q$, and $N\le6$, the largest certified norm bound is
$7.717202888999335\times10^{-13}$.
\end{description}

The following statements remain open:

\begin{enumerate}
\item a useful estimate uniform as $N\to\infty$;
\item positivity of every infinite-cutoff finite section;
\item continuity and closure of the complete arithmetic form on the required
domain;
\item the implication from its null structure to the canonical native
PhaseNav closure;
\item the complete bridge represented by \texttt{SOH-C005}.
\end{enumerate}

Accordingly, controlled cutoff removal is now available for each fixed finite
section, but this chapter does not establish global Weil positivity and does
not prove the Riemann Hypothesis.
